
\documentclass[12pt]{book}
%%%%%%%%%%%%%%%%%%%%%%%%%%%%%%%%%%%%%%%%%%%%%%%%%%%%%%%%%%%%%%%%%%%%%%%%%%%%%%%%%%%%%%%%%%%%%%%%%%%%%%%%%%%%%%%%%%%%%%%%%%%%
%\usepackage{sw20etpl}

%TCIDATA{TCIstyle=book/book.lat,etpl,sw20etp}

%TCIDATA{Created=Sat Aug 30 22:02:35 1997}
%TCIDATA{LastRevised=Mon Sep 01 00:05:13 1997}

%\input{tcilatex}
\begin{document}


\chapter{Review on Hypothesis Testing}

\section{Some Important Concepts}

\begin{description}
\item[Null Hypothesis $\left( H_{0}\right) $]  A maintained hypothesis that
is held to be true unless sufficient evidence to the contrary is presented.
This maintained claim usually carries some form of \emph{significance} or 
\emph{importance}.

\item[Alternative Hypothesis $\left( H_{1}\right) $]  A hypothesis that is
held to be true when the null hypothesis is rejected.

\item[Simple Hypothesis]  A hypothesis that specifies a single value for a
population parameter of interest.

\item[Composite Hypothesis]  A hypothesis that specifices a range of values
for a population parameter of interest.

\item[One-sided Alternative]  An alternative hypothesis involving all
possible values of a population parameter on either one side or the other of
the value specified by a simple null hypothesis.

\item[Two-sided Alternative]  An alternative hypothesis involving all
possible values of a population parameter other than the value specified by
a simple null hypothesis.

\item[Type I Error]  The mistake of rejecting a true null hypothesis.

\item[Type II Error]  The mistake of accepting a false null hypothesis.

\item[Significance Level $\left( \alpha \right) $]  The probability of
rejecting a true null hypothesis. This is usually kept at a lower value
closer to zero.

\item[Power $\left( 1-\beta \right) $]  The probability of rejecting a false
null hypothesis. This is determined by both the chosen $\alpha $ level and
the testing procedure.

\item[$p$-value / probability value]  The smallest significance level at
which a null hypothesis can be rejected. Alternative, one can think of it as
the level of significance $\alpha $ when the evidence is used as a critical
value.
\end{description}

\section{The 2x2 Decision Paradigm}

\begin{center}
\begin{tabular}{|ll|c|c|}
\hline
& \multicolumn{1}{l|}{} & \multicolumn{2}{|c|}{True State of Nature} \\ 
\cline{3-4}\cline{4-4}
& \multicolumn{1}{l|}{} & $H_{0}\,$true & $H_{0}$ false \\ \hline
\multicolumn{1}{|c}{Decision} & \multicolumn{1}{|c|}{
\begin{tabular}{l}
Accept $H_{0}$ \\ 
Reject $H_{0}$%
\end{tabular}
} & 
\begin{tabular}{l}
\begin{tabular}{l}
Correct \\ 
Probability $=\left( 1-\alpha \right) $%
\end{tabular}
\\ 
\begin{tabular}{l}
Wrong (type I error) \\ 
Level of significance $=\alpha $%
\end{tabular}
\end{tabular}
& 
\begin{tabular}{l}
\begin{tabular}{l}
Wrong (type II\ error) \\ 
Probability $=\beta $%
\end{tabular}
\\ 
\begin{tabular}{l}
Correct \\ 
Power $=1-\beta $%
\end{tabular}
\end{tabular}
\\ \hline
\end{tabular}
\end{center}

\section{Various Types of $H_{0}$ and $H_{1}$}

\begin{description}
\item[Case I]  A simple null vs. a one-sided upper-tailed composit
alternative:

$H_{0}:\theta =\theta _{0};\qquad H_{1}:\theta >\theta _{0}$

\item[Case II]  A simple null vs. a one-sided lower-tailed composit
alternative:

$H_{0}:\theta =\theta _{0};\qquad H_{1}:\theta <\theta _{0}$

\textbf{Remark: }Testing procedure is equivalent to Case I.

\item[Case III]  A composit null vs. a one-sided upper-tailed composit
alternative:

$H_{0}:\theta \leq \theta _{0};\qquad H_{1}:\theta >\theta _{0}$

\item  \textbf{Remark: }Testing procedure is exact opposite of Case I.

\item[Case IV]  A composit null vs. a one-sided lower-tailed composit
alternative:

$H_{0}:\theta \geq \theta _{0};\qquad H_{1}:\theta <\theta _{0}$

\item  \textbf{Remark: }Testing procedure is equivalent to Case III or exact
opposite of Case I.

\item[Case V]  A simple null vs. a two-sided composit alternative:

$H_{0}:\theta =\theta _{0};\qquad H_{1}:\theta \neq \theta _{0}$

\textbf{Remark: }Testing procedure is the combination of Case I and III.
\end{description}

\section{General Hypothesis Testing Procedure}

\begin{description}
\item[Step 1:]  Write down $H_{0}$ and $H_{1}$ (i.e. express the claim in
terms of the population paramter(s)).

\item[Step 2:]  Identify the sample and test statistics (i.e. the evidence)
to be used.

\item[Step 3:]  Determine the sampling distribution of the sample and/or
test statistics under $H_{0}$ (i.e. plot the distribution assuming $H_{0}$
is true). Block out the acceptance and rejection regions in the domain of
the sampling distribution.

\item[Step 4:]  If $\alpha \,$is given, determine the critical value(s) in
terms of the sample and/or test statistics. If $\alpha $ is not provided,
find the $p$-value using the evidence as a critical value.

\item[Step 5]  Make decision and draw your conclusion in the context of the
problem.
\end{description}

\end{document}
